\documentclass[11pt]{article}
\usepackage[utf8]{inputenc}
\usepackage[T1]{fontenc}
\usepackage{lmodern}
\usepackage[margin=1in]{geometry}
\usepackage{amsmath,amssymb}
\usepackage{textcomp}
\usepackage[dvipsnames]{xcolor}
\usepackage{enumitem}
\usepackage{parskip}
\usepackage{titlesec}
\usepackage[colorlinks=true,linkcolor=blue,citecolor=blue,urlcolor=blue]{hyperref}

\titleformat{\section}{\large\bfseries}{\thesection}{1em}{}
\titleformat{\subsection}{\normalsize\bfseries}{\thesubsection}{1em}{}
\titleformat{\subsubsection}{\normalsize\bfseries\itshape}{}{0em}{}

\title{\textbf{Survive Every Time, Break Once:\\ Measuring the Adversarial Boundary and Cost of Memory-Pinning under Context Compaction}}
\author{}
\date{}

\newcommand{\ph}{\textit{[placeholder --- fill in after experiments are run]}}

\begin{document}
\maketitle
\vspace{-2em}

\begin{quote}
\small\textit{Assembled from the six confirmed stage outputs of an \texttt{auto-research-discussion} run (seed topic T3.1), post research-taste audit (\textcolor{ForestGreen}{\textbf{PASS}}). Status: \textbf{specification / paper draft --- no experiments have been run.} All numbers below are pre-registered design targets, not results.}
\end{quote}

\begin{abstract}
Long-horizon LLM agents periodically compact their conversation and tool history to stay within a token budget, and memory-pinning is the mechanism operators lean on to keep a safety- or task-critical constraint alive through every such compaction pass. We show this reliance rests on an unexamined asymmetry: under a fixed-budget lossy compaction channel replayed across a session, a pinned constraint survives only if it re-wins the retention-budget competition at every compaction event, whereas an adversary needs it to lose only once, and can amortize an offline optimization budget to engineer that single loss against a defender who pays only a bounded per-event budget. Governance Decay (arXiv 2606.22528) --- the work that introduced compaction-induced constraint erasure, the compaction-eviction attack, and the constraint-pinning defense whose components our variants decompose --- already shows an optimization-based injection can defeat pinning on every tested model, but not where along an attack-strength axis that break happens, nor what it costs the attacker relative to what the defense costs the defender to run. We build a measurement harness that sweeps a normalized injection-strength parameter across five pinning-style defense variants and two attack families (a gradient-optimized suffix search and a template-role-forgery attack), verify constraint violation with a two-tier oracle gated on the constraint's state before and after compaction, fit a break-through-rate boundary per variant with bootstrap confidence intervals, and convert attacker and defender effort into a common cost currency to report the resulting cost ratio at that boundary. Budget-matched benign-filler and no-defense controls isolate the adversarial cause from generic context-length degradation. We report where, and how cheaply, each pinning variant fails, and what this implies for treating pinning as a sufficient safeguard for autonomous agent deployments.
\end{abstract}

\section{Introduction}

Long-horizon LLM agents now routinely run sessions long enough that their conversation and tool-call history must be periodically compacted --- summarized and compressed --- to stay within a usable context budget. Because that compaction is lossy by design, a growing set of deployments rely on memory-pinning: re-inserting, redundantly copying, checksum-verifying, or structurally isolating a small set of safety- or task-critical constraints so they survive every compaction pass rather than being silently summarized away. The appeal of pinning is that it is cheap and legible --- a fixed rule like ``never call this tool'' or ``always defer above a spending threshold'' can be attached once and, the assumption goes, will hold for the rest of the session.

Yet a load-bearing safety assumption behind this design remains unmapped. Governance Decay (arXiv 2606.22528) established the phenomenon this paper studies: it introduced compaction-induced governance decay, demonstrated the compaction-eviction attack (adversarial content that biases the summarizer into dropping a policy), proposed constraint pinning as the defense --- a pinned buffer exempt from compaction, re-injected verbatim, with per-step integrity checking, from which this paper's re-injection, checksum-verified, and structural-slot variants are directly decomposed --- reported its defender-side cost ($\approx$47 tokens, under 0.5\% overhead), and showed that an optimization-based adversarial injection defeats pinning on every tested model, with operator impersonation in recent context as the residual failure mode. So pinning is known not to be unconditionally robust, and its defense taxonomy and defender cost are inherited from that work, not contributed here. What Governance Decay does not establish is where, along a continuum of attacker effort, the break happens for a given pinning mechanism, or what it costs the attacker to reach that point relative to what the defense costs the defender to run every turn: its headline is a binary ``pinning works at negligible cost (except under impersonation)'' at fixed attack configurations. This gap matters because a defender choosing among pinning variants is not choosing between ``robust'' and ``not robust'' --- every variant eventually fails --- but between variants that differ in how much attacker effort a break requires and how much overhead the defense itself imposes to buy that resistance. The precise shape of that trade-off, not merely the fact that pinning eventually fails, is the safety-relevant quantity nobody has measured.

The nearest existing work approaches this gap from adjacent but distinct angles. Governance Decay, as just discussed, is the nearest neighbor overall --- it owns the phenomenon, the attack premise, the pinning defense, and the defender-cost figure, but reports a binary defended/undefended outcome rather than a break boundary on a continuous strength axis or an attacker/defender cost ratio. IH-Benchmark (arXiv 2607.25987) is the nearest oracle instrument --- a predicate-DSL check plus category-scoped LLM judges for instruction-hierarchy conflicts --- but it scores single-decision conflict outcomes, not persistence of a constraint across a compaction event. Security--Fidelity Tradeoffs (arXiv 2606.30783) prices the fidelity/utility cost of prompt-injection defenses, the defender half of a cost account, without an attacker-side denominator. TensorTrust builds an interpretable benchmark from a competitive human-authored attack/defense game and explicitly measures a robustness-vs-usability trade-off --- the closest existing evaluation design to what this paper needs --- but its defenses are static, human-authored, and single-turn, with no compaction dynamic and no continuous attack-strength axis to sweep. LightMem is the closest existing mechanism: its ``sleep-time'' consolidation step is structurally similar to a protected-memory slot that this paper's structural-slot variant is modeled on, but LightMem is evaluated only on benign long-horizon dialogue consistency, never under adversarial stress. Optimization-based jailbreak attacks such as SlotGCG and TAO-Attack, and the canonical GCG suffix-search method they extend, optimize against single-turn safety-refusal at a fixed attack budget rather than sweeping strength against a structured pinning defense across compaction events. Template-role-forgery attacks such as ChatInject demonstrate that forged chat-role structure breaks agent defenses at one fixed configuration, but neither sweep a strength parameter nor target memory-pinning specifically. And a single-backbone-call abstraction like ThreatSnapshots deliberately isolates one LLM call to fix attribution, which cannot represent a pinned constraint's survival --- a property of a repeated, multi-event lifecycle, not a single call.

Turning this gap into a measurable question surfaces its own difficulties. Compaction is lossy summarization by construction, so any pinning mechanism either competes for retention budget inside that lossy pass or must be exempted from it entirely --- and no metric distinguishes which strategy a given variant is actually using without instrumenting per-turn constraint state directly. The two attack families that can defeat pinning are structurally incomparable: a gradient-optimized suffix search is measured in optimization iterations, candidate-set size, and token budget, while a template-role-forgery attack is measured in forged-turn count and payload size, with no natural shared unit between them. Verifying that a constraint was actually violated after compaction requires an oracle that checks a relational property across two points in time --- present before compaction, violated at or after it; two-tier oracles combining deterministic checks with scoped LLM judges already exist (IH-Benchmark is the nearest), but they verify single-decision outcomes, not this before/after-compaction persistence property. Attacker and defender effort are denominated in entirely different resources (offline optimization compute versus per-turn runtime overhead), so no currency for comparing them is uniquely correct. And because LLM outputs and any LLM-judge oracle are non-deterministic, a break-through rate is a stochastic estimand whose apparent boundary can be an artifact of an underpowered trial count rather than a real property of the mechanism being tested.

Our approach turns each of these obstacles into a specific design choice rather than a hedge. Five pinning variants --- a no-defense control, verbatim re-injection, redundant $k$-copy pinning, checksum-verified pinning, and a structurally isolated protected-memory slot --- are implemented as interchangeable middleware, instrumented to log each constraint's presence and verbatim state before and after every compaction event, so which strategy a variant is actually using is directly observable rather than inferred. The two attack families' raw, incomparable strength units are normalized onto one shared log-scaled axis, so a single break-through-vs-strength curve can be fit and compared per variant regardless of which family produced it. Constraint violation is verified by a two-tier oracle --- a deterministic check for hard constraints and a rubric-based judge for soft constraints --- gated by an explicit before/after-compaction attribution requirement, so a persistence claim is checked as the relational property it actually is. Attacker effort (to reach a variant's boundary) and defender effort (imposed every turn) are converted into a token-equivalent currency and, separately, an API-dollar-class currency, so the ratio between them is a single reportable number whose currency-sensitivity is itself checked rather than assumed away. And every boundary is fit only after a power gate raises the per-cell trial count until sampling noise is small enough that the fitted boundary is distinguishable from a flat rate with the same shape.

We evaluate this design across tool-using agent benchmarks with explicit policy constraints and long multi-session dialogue benchmarks that force repeated compaction events, comparing the adversarial injection against budget-matched benign-filler and no-defense-under-benign-compaction controls that isolate the adversarial cause from generic context-length degradation.

In summary, we contribute:
\begin{itemize}[leftmargin=1.4em]
\item An identification and formalization of the durability/cost asymmetry underlying memory-pinning as a safeguard: a pinned constraint must re-win a retention competition at every compaction event, while an adversary needs it to lose only once and can amortize an offline attack budget against a bounded per-event defense budget.
\item A strength-normalization scheme that puts two structurally incomparable attack families --- gradient-optimized suffix search and template-role forgery --- on one comparable axis, letting a single break-through boundary be fit and compared per pinning variant regardless of attack family.
\item An attribution-gated persistence oracle for constraints across a compaction event. The two-tier structure itself (deterministic hard-constraint check + rubric-based judge) follows existing instruments --- IH-Benchmark's predicate DSL plus category-scoped judges is the nearest --- and the contribution is confined to the relational before/after-compaction gate that turns a single-decision check into a persistence check.
\item A dual-currency cost-accounting method yielding an attacker/defender cost ratio at each variant's measured boundary. Defender-side cost alone is already reported by Governance Decay, and defense-cost pricing exists in Security--Fidelity Tradeoffs; the contribution is the attacker-to-defender ratio at a measured break boundary, together with a check of its sensitivity to the chosen currency.
\item An empirical account of where along the attack-strength axis, and at what cost, each pinning variant fails under adversarial context compaction --- noting that three of the five variants (re-injection, checksum-verified, structural slot) decompose Governance Decay's own proposed defense, so the finding prices an existing defense rather than a new one --- and what this implies for treating pinning as a sufficient safeguard.
\end{itemize}

\section{Problem Definition}

\noindent\textbf{Formal setup.} Consider an agent trajectory of turns $\tau_1,\dots,\tau_T$ with a compaction event $\kappa$ that maps history to at most $B$ tokens, occurring $m \geq 1$ times over the trajectory. A pinned constraint $c$, drawn from a suite $C$, must hold for the remainder of the session once introduced. A defense variant $d$ is drawn from $D = \{d_0, d_{\text{reinj}}, d_{\text{kcopy}}, d_{\text{cksum}}, d_{\text{slot}}\}$ (no-defense control, verbatim re-injection, redundant $k$-copy pinning, checksum/tag-verified pinning, structurally isolated protected-memory slot). An adversarial injection procedure, drawn from attack family $a \in A = \{a_{\text{wb}}, a_{\text{bb}}\}$ (white-box gradient-optimized suffix search; black-box template-role forgery), is parametrized by a normalized strength scalar $s \in [0,1]$. A break-through indicator $V(c,T;d,a,s,\omega) \in \{0,1\}$ records whether $c$ was violated at or after a compaction event, given it was present immediately before that event, for trial seed $\omega$. The object of study is, for each (variant, attack family, model $g$), the break-through-rate function $p_{d,a,g}(s) = \Pr[V=1]$: does it cross a meaningful threshold within the tested strength range, and if so, at what strength (the boundary, BT50/BT90), and at what ratio $R$ of attacker to defender cost at that boundary?

\noindent\textbf{Why this is hard.} Compaction is lossy summarization by construction, so a defense variant either competes for retention budget inside that lossy pass ($d_{\text{reinj}}$, $d_{\text{kcopy}}$, $d_{\text{cksum}}$ all attempt post-hoc recovery after competing) or is exempted from it entirely ($d_{\text{slot}}$) --- and no metric distinguishes which of these two strategies is responsible for a variant's observed resistance without directly instrumenting per-turn constraint state. The two attack families are measured in structurally incomparable raw units --- optimization iterations $\times$ candidate-set size $\times$ token budget for the gradient family, versus forged-turn count $\times$ payload token budget for the template family --- so a boundary reported in either family's raw units is not comparable to the other's without an explicit, disclosed normalization. Verifying $V$ requires an oracle that checks a relational property across two time points (present before $\kappa$, violated at or after it), which is a materially different object than a single-turn harmfulness judgment or a single-backbone-call security score. The attacker's effort (offline optimization compute, queries, or tokens to reach a given $s$) and the defender's effort (extra tokens and latency imposed every turn by running variant $d$) are denominated in different kinds of resources entirely, so no currency for their ratio is uniquely correct --- any chosen currency embeds a methodological assumption that must be surfaced and stress-tested, not hidden. And because $p_{d,a,g}(s)$ is estimated from a finite, noisy set of trials, an apparent boundary can be an artifact of an underpowered trial count rather than a real property of the (variant, attack, model) cell, unless the trial count is set by an explicit power criterion rather than fixed in advance.

\section{Method Overview}

At a conceptual level, the study is a single measurement pipeline with five cooperating parts. A matrix of pinning-style defense variants and a normalized cross-family attack-strength axis together produce, per condition, a sequence of trial outcomes; a two-tier, attribution-gated oracle turns each trial into a verified break-through indicator; a power-gated, bootstrap-interval curve fit turns the resulting sequence into a comparable boundary scalar per variant; and a dual-currency cost accounting turns attacker and defender effort into a single reportable ratio at that boundary.

The \textbf{structural defense matrix} implements the five pinning variants as interchangeable middleware around a compaction summarizer, instrumented to log each constraint's presence and verbatim state before and after every compaction event, so which recovery strategy a variant is actually using is directly observable. The \textbf{strength-normalization scheme} takes each attack family's raw, structurally incomparable strength unit and maps it onto one shared log-scaled $[0,1]$ axis, so a break-through curve fit under one family is directly comparable to a curve fit under the other. The \textbf{two-tier violation oracle} takes each trial's post-compaction state and verifies whether the pinned constraint was violated, using a deterministic check for hard constraints and a rubric-based judge for soft constraints, gated by an explicit before/after-compaction attribution requirement. The \textbf{dual-currency cost accounting} takes the attacker effort needed to reach a variant's boundary and the defender overhead that variant imposes every turn, converts both into a token-equivalent currency and, separately, an API-dollar-class currency, and reports the ratio under both. The \textbf{power-gated boundary estimation} takes the resulting break-through series and fits a logistic curve per (variant, model, attack family), raising the per-cell trial count until sampling noise is small enough for the fitted boundary to be trustworthy, then reports it with a bootstrap confidence interval.

\section{Methodology}

\subsection{Structural Defense Matrix}
\textbf{Challenge.} Compaction is lossy summarization by construction, so any pinning mechanism either competes for retention budget inside that lossy pass or must be exempted from it entirely, and no existing metric distinguishes which strategy a given variant is actually using.

\textbf{Insight.} Only a variant that removes pinned content from the compaction pass entirely sidesteps the lossy competition; every other variant competes inside it and differs only in how it attempts post-hoc recovery --- a distinction invisible unless per-turn constraint state is instrumented directly rather than inferred from a final transcript.

\textbf{Method.} Five pinning variants are implemented as interchangeable middleware wrappers around a compaction summarizer: a no-defense control, verbatim re-injection of the pinned constraint after each compaction event, redundant $k$-copy pinning, checksum/tag-verified pinning (a validator recomputes a hash of the protected span and re-inserts on mismatch), and a structurally isolated protected-memory slot exempt from summarization. An agent harness logs, per turn and per constraint, the constraint's presence and verbatim state immediately before and after every compaction event.

\subsection{Strength-Normalization Scheme}
\textbf{Challenge.} The two attack families capable of defeating pinning are measured in structurally incomparable raw units, so a boundary reported in one family's units cannot be directly compared to a boundary reported in the other's.

\textbf{Insight.} Normalizing two structurally incomparable raw attack budgets onto one log-scaled axis is what allows a single break-through-vs-strength curve to be fit and compared per variant regardless of which attack family produced it; reporting each family's raw units instead would leave every variant incomparable across families.

\textbf{Method.} Each attack family is parametrized with an explicit strength knob --- optimization iterations $\times$ candidate-set size $\times$ token budget for the gradient-optimized suffix family; forged-role/turn count $\times$ payload token budget for the template-role-forgery family --- and both are mapped onto a shared $[0,1]$ axis via a log-transformed, per-family-bounded scale: $s = (\log \rho_a - \log \rho_a^{\min}) / (\log \rho_a^{\max} - \log \rho_a^{\min})$, where $\rho_a$ is the raw attack budget for family $a$ and the bounds are fixed before the sweep is run.

\subsection{Two-Tier Violation Oracle}
\textbf{Challenge.} Verifying that a pinned constraint was violated after compaction requires checking a relational property across two points in time --- present before the compaction event, violated at or after it --- which is a materially different object than a single-turn harmfulness judgment or a single-backbone-call security score.

\textbf{Insight.} A persistence oracle checks a relational, two-time-point property. The two-tier structure --- deterministic checks plus scoped LLM judges --- is not new (IH-Benchmark's predicate DSL with category-scoped judges is the nearest existing instrument); what no existing oracle supplies is the relational gate that conditions the verdict on the constraint being present immediately before the compaction event and violated at or after it.

\textbf{Method.} Constraint violation is verified with a two-tier oracle: a deterministic programmatic check for hard constraints (a forbidden tool call fired, a numeric or format rule broken) and a rubric-based LLM-judge checklist, adapted from guideline-based jailbreak-evaluation methodology, for soft constraints. Both tiers are gated by an attribution requirement that the constraint was present immediately before the compaction event and violated at or after it, and the judge tier is audited against a held-out set of human-labeled cases before being used for confirmatory claims.

\subsection{Dual-Currency Cost Accounting}
\textbf{Challenge.} Attacker effort (offline optimization compute) and defender effort (per-turn runtime overhead) are denominated in entirely different kinds of resources, so reporting them separately leaves ``is this variant worth its overhead'' unanswerable as a single number, and no currency for combining them is uniquely correct. Each side in isolation is already priced by prior work --- Governance Decay reports pinning's defender cost ($\approx$47 tokens, $<$0.5\% overhead) and Security--Fidelity Tradeoffs prices the fidelity cost of injection defenses --- but no prior work relates attacker cost to defender cost at a measured break point.

\textbf{Insight.} Expressing attacker and defender effort in the same unit, and doing so under two independently motivated currencies rather than one, turns an incommensurable comparison into a checkable number whose currency-sensitivity is reported rather than hidden.

\textbf{Method.} Attacker cost is computed as the optimizer queries, forward/backward passes, or API tokens spent to reach a variant's measured boundary; defender cost is computed as the extra tokens and added latency that variant imposes every turn. Both are converted into a token-equivalent common currency and, separately, into an API-dollar-class currency, and the attacker-to-defender ratio $R$ is computed at each variant's boundary under both currencies.

\subsection{Power-Gated Boundary Estimation}
\textbf{Challenge.} A break-through rate is a stochastic estimand, so an apparent boundary can be an artifact of an underpowered trial count rather than a real property of the (variant, attack, model) cell being tested.

\textbf{Insight.} Reporting a boundary as a bootstrap interval derived from a trial count set by an explicit power criterion, rather than a point estimate from an arbitrarily chosen number of trials, is what lets a genuine boundary be distinguished from an equally-shaped curve that is actually sampling noise around a flat rate.

\textbf{Method.} For each (variant, model, attack family), the per-cell trial count is raised until the Wilson score interval's half-width falls below a preset threshold at the steepest point of the observed grid, and only then is a logistic break-through-vs-strength curve fit, yielding a boundary scalar (BT50/BT90) with a bootstrap confidence interval, alongside the area under the break-through curve as a strength-agnostic summary.

\section{Experiments}

\subsection{Experimental Setup}

\textbf{Datasets.} AgentDojo --- a tool-using agent environment that natively hosts task constraints alongside a security-relevant execution log, giving the hard-constraint oracle tier a programmatically checkable substrate without relying on judge-only scoring. InjecAgent --- an indirect-prompt-injection agent benchmark supplying the injection substrate and attack-insertion scaffolding the template-role-forgery attack family is adapted from. $\tau$-bench --- tool-agent tasks with explicit policy constraints closer to production agent workloads, testing whether the boundary and cost findings generalize beyond one benchmark's task distribution. LongMemEval and LoCoMo --- long multi-session dialogue benchmarks that natively require multiple, repeated compaction events, which single-session benchmarks do not force; the cross-condition ranking-stability claim specifically needs a workload where a pinned constraint must re-survive more than one compaction event, since a single-compaction task cannot distinguish ``survived once'' from ``reliably re-wins the retention competition every time.''

\textbf{Baselines and controls} (sanity checks and confound-isolation controls, not competitors to beat --- this is a measurement study, and no research question here is framed as ``our procedure beats baseline X''). A no-defense variant under benign, non-adversarial compaction (no attacker present) establishes what fraction of constraints are lost to compaction alone; this is the necessary control for the framing-tier attribution question below. A budget-matched benign-filler injection replaces the attacker's payload with benign filler of identical token budget, count, and insertion position, isolating whether break-through is caused by the adversarial content specifically or by generic long-context degradation at matched budget. Scripted, non-adaptive template variants (a fixed forged-role count, no optimization) check that any observed strength-dependent trend is not simply an artifact of longer prompts in general.

\textbf{Attack-arm scope.} The gradient-optimized suffix family requires gradient access and is therefore confined to open-weight backbones; closed-weight API models (including the $\tau$-bench-style production workloads) are covered only by the black-box template-role-forgery family. RQ5's cross-model ordering claim is accordingly family-conditioned: for the gradient family it is a claim about open-weight models only, and generalization to API models rests on the template family alone.

\textbf{Metrics.} The two-tier-verified break-through rate with Wilson 95\% confidence intervals; the boundary scalars BT50 and BT90 with bootstrap confidence intervals; the area under the break-through curve; the attacker/defender cost ratio at BT50 under two currencies; Cohen's $\kappa$ for oracle judge-human agreement; and Kendall's $\tau$ for cross-model/cross-attack-family ordering stability.

\textbf{Research questions:}
\begin{enumerate}[leftmargin=1.6em]
\item RQ0 (framing) --- Is observed post-compaction break-through attributable to the adversarial injection specifically, rather than to generic benign compaction-induced degradation?
\item RQ1 --- Does the resistance ordering across the five pinning variants follow the predicted structural-slot-highest ranking?
\item RQ2 --- Is a reliable break-through boundary fittable for most (variant, attack-family, model) cells under the normalized strength axis?
\item RQ3 --- Does the two-tier oracle's judge tier agree with human judgment closely enough to support confirmatory claims?
\item RQ4 --- Does the attacker/defender cost ratio differ meaningfully across variants, and does the ranking hold across cost currencies?
\item RQ5 --- Is the cross-model, cross-attack-family ordering of variants by boundary and cost ratio stable?
\end{enumerate}

\subsubsection{RQ0: Is observed break-through attributable to the adversarial injection specifically, rather than to generic benign compaction-induced degradation?}
\textbf{Setup.} Compare the no-defense variant's break-through rate under benign (non-adversarial, empty-payload) compaction against its break-through rate under the adversarial injection at matched strength and turn count; separately compare the adversarial injection's break-through rate against the budget-matched benign-filler control at identical token budget, count, and insertion position.

\textbf{Expected outcome.} The benign-compaction break-through rate stays low (its Wilson upper bound remains below 0.20), and the adversarial injection's break-through rate is significantly higher than the budget-matched benign-filler control's at matched budget --- establishing that break-through is attributable to the adversarial content, not to context degradation or compaction alone.

\textbf{Falsification condition.} Refuted if the benign-compaction break-through rate's Wilson upper bound reaches or exceeds 0.20 (constraints are already lost to compaction alone, making the adversarial framing moot), or if the adversarial injection's break-through rate is statistically indistinguishable from the budget-matched benign-filler control's (overlapping Wilson confidence intervals) --- either would mean the claimed adversarial cause is not isolated from its strongest alternative explanation.

\textbf{Hypothesis (plain language).} Neither of the two benign-compression views in the surrounding literature models an adversary trying to evict a protected constraint, so their retention guarantees say nothing about adversarial durability; this RQ verifies directly that what is measured here is adversarial break-through, not an artifact those benign models would already predict.

\textbf{Metrics.} Break-through rate under benign compaction (Wilson CI); break-through rate under budget-matched benign filler vs. adversarial injection at matched budget (Wilson CIs, overlap check).

\textbf{Result.} \ph

\subsubsection{RQ1: Does the resistance ordering across the five pinning variants follow the predicted structural-slot-highest ranking?}
\textbf{Setup.} Fit BT50 for all five variants per (model, attack family) cell using the power-gated boundary estimation method; compare the five BT50 values pairwise within each cell.

\textbf{Expected outcome.} The five variants' BT50 values are ordered $d_{\text{slot}} > d_{\text{cksum}} > d_{\text{kcopy}} > d_{\text{reinj}} > d_0$ in a majority of tested cells --- resistance increases with how far a variant's recovery mechanism is removed from relying on the lossy compaction pass to have preserved the constraint verbatim.

\textbf{Falsification condition.} Refuted if the predicted ordering does not hold in a majority of tested cells, or if any adjacent pair in the predicted order is reversed with a non-overlapping bootstrap confidence interval in a majority of cells --- either would mean resistance does not track how far a variant's recovery mechanism sits from the lossy pass, and the structural-exemption insight does not hold as stated.

\textbf{Hypothesis (plain language).} Resistance to adversarial break-through is not uniform across pinning implementations --- the way a variant relates to the lossy compaction pass (exempted from it versus competing inside it) predicts its relative robustness.

\textbf{Metrics.} Per-variant BT50 with bootstrap CI; pairwise ordering agreement rate across cells.

\textbf{Result.} \ph

\subsubsection{RQ2: Is a reliable break-through boundary fittable for most (variant, attack-family, model) cells under the normalized strength axis?}
\textbf{Setup.} For every tested (variant, attack-family, model) cell, fit the logistic break-through-vs-strength curve on the normalized strength axis and evaluate the fit-quality bar (McFadden pseudo-$R^2 \geq 0.3$, significantly positive slope, bounded BT50 confidence interval).

\textbf{Expected outcome.} At least 4 of the 5 defense variants (per model, per attack family) clear the fit-quality bar, with McFadden pseudo-$R^2 \geq 0.3$, a significantly positive slope (Holm-corrected $p<0.05$), and a BT50 bootstrap 95\% confidence interval no wider than 25\% of the tested strength range.

\textbf{Falsification condition.} Refuted if fewer than 4 of the 5 variants clear the fit-quality bar in a majority of (model, attack-family) cells; a variant whose break-through rate instead stays flat and below a low ceiling across the entire tested strength range is not a falsifier of this RQ but a distinct, equally reportable unbreakable-in-range finding, to be reported separately rather than counted as a failed fit.

\textbf{Hypothesis (plain language).} The normalized strength axis is comparable enough across attack families that a genuine, well-fit boundary --- not an artifact of the normalization choice --- is findable for most pinning variants.

\textbf{Metrics.} McFadden pseudo-$R^2$ per cell; slope significance; BT50 bootstrap CI width as a fraction of the strength range.

\textbf{Result.} \ph

\subsubsection{RQ3: Does the two-tier oracle's judge tier agree with human judgment closely enough to support confirmatory claims?}
\textbf{Setup.} Ground the judge's checklist categories in the pinned constraint's own written specification, draft the rubric on a $\sim$50-case held-out calibration slice, then audit the frozen rubric against $\geq 2$ independent human raters on a disjoint $\sim$200-case held-out audit set, computing Cohen's $\kappa$ between judge and human labels and inspecting the confusion structure for a systematic collapse onto one label.

\textbf{Expected outcome.} The judge tier achieves Cohen's $\kappa \geq 0.7$ against human labels on the held-out audit set (an instrument-validation sequence raised this bar from an initial 0.6 to the 0.7 minimum standard for confirmatory inter-rater agreement), with no systematic collapse in the disagreement structure.

\textbf{Falsification condition.} Refuted if Cohen's $\kappa$ falls below 0.7, or if disagreements show a systematic collapse onto one label even with $\kappa \geq 0.7$; in either case, soft-constraint results are demoted to secondary/exploratory status and only hard-constraint (programmatically verified) results are treated as confirmatory for RQ1, RQ2, RQ4, and RQ5.

\textbf{Hypothesis (plain language).} A persistence oracle built from a deterministic hard-constraint check plus a rubric-based judge for soft constraints is reliable enough to support the study's confirmatory claims, not merely its exploratory ones.

\textbf{Metrics.} Cohen's $\kappa$ (judge vs. human); per-category disagreement rate.

\textbf{Result.} \ph

\subsubsection{RQ4: Does the attacker/defender cost ratio differ meaningfully across variants, and does the ranking hold across cost currencies?}
\textbf{Setup.} Compute the attacker/defender cost ratio $R$ at each variant's measured BT50 under both the token-equivalent and API-dollar-class currencies; compare ratios pairwise across variants with a Holm-corrected Wilcoxon signed-rank test on $\log R$, across tasks.

\textbf{Expected outcome.} The cost ratio differs by at least $2\times$ between the cheapest-to-break and most-expensive-to-break variant (Holm-corrected Wilcoxon rejecting equality for at least one pair, $\alpha = 0.05$), and the resulting variant ranking is identical (Kendall $\tau = 1$) across the two currencies.

\textbf{Falsification condition.} Refuted if the cost ratio is statistically indistinguishable across all variant pairs (no pair reaches the $2\times$ threshold with a significant test), which would mean pinning-mechanism choice does not robustly change adversarial cost; separately refuted (as a ranking-stability failure) if the variant ranking reverses between the two currencies, in which case a single ratio is not reported and currency-conditional rankings are reported side by side instead.

\textbf{Hypothesis (plain language).} Pinning implementations differ meaningfully in cost-effectiveness even though every variant eventually fails --- some are more expensive to break, per unit of defender overhead, than others.

\textbf{Metrics.} Cost ratio per variant under each currency with bootstrap CI; Wilcoxon test statistic and $p$-value; Kendall $\tau$ between currency-conditional rankings.

\textbf{Result.} \ph

\subsubsection{RQ5: Is the cross-model, cross-attack-family ordering of variants by boundary and cost ratio stable?}
\textbf{Setup.} Compute the variant ordering by BT50 and, separately, by cost ratio, for each tested model and attack family; compute the mean pairwise Kendall's $\tau$ of these orderings across all model pairs and both attack families, with a permutation test for significance.

\textbf{Expected outcome.} Mean pairwise Kendall's $\tau \geq 0.6$ at permutation $p<0.05$ for both the BT50-based and cost-ratio-based orderings, indicating the relative robustness ranking of pinning variants is largely model- and attack-family-agnostic even as absolute BT50 values shift.

\textbf{Falsification condition.} Refuted if mean pairwise Kendall's $\tau$ is not significantly positive ($\tau \approx 0$ or negative, or fails the permutation test) --- meaning which pinning variant is best depends on the specific model and attacker, and no transferable deployment guidance can be given from a single ranking.

\textbf{Hypothesis (plain language).} The relative robustness ranking of pinning variants transfers across models and attack families, even though absolute break points shift --- giving practitioners guidance that does not require re-running the full sweep for every new backbone.

\textbf{Metrics.} Mean pairwise Kendall's $\tau$ (BT50-based and cost-ratio-based) with permutation $p$-value.

\textbf{Result.} \ph

\vspace{0.5em}
\noindent\textbf{Further evaluation.} With more time, a natural extension is to calibrate the gap between warm-started/transferable gradient-attack suffixes (used here to keep the white-box arm tractable) and a fully re-optimized per-cell attack, since amortization could understate attack strength and bias BT50 upward; a representative-subset comparison of warm-started vs. fully-re-optimized attacks would quantify and correct for this offset rather than assume it away.

\section{Related Work}

\noindent\textit{Citation-provenance caveat (partially discharged 2026-08-21).} Five records in this section are now arXiv-API-verified and added to this topic's reference cache: Governance Decay (2606.22528), IH-Benchmark (2607.25987), Security--Fidelity Tradeoffs (2606.30783), ChatInject (2509.22830), and SlotGCG (2606.05609). The remaining prior-work names --- ThreatSnapshots, TAO-Attack, GCG, GuidedBench, TensorTrust, COMI, the companion information-channel view, and LightMem --- were grounded during ideation via a general literature corpus separate from this topic's formal reference cache and remain provisional attack/defense \emph{mechanism} references (gradient suffix search; template-role forgery; benign compaction/memory pipelines) pending a targeted re-harvest; none is treated as a numerical baseline this study competes against.

The closest prior work to this paper is Governance Decay (arXiv 2606.22528), and the delta must be stated plainly before surveying the wider landscape. Governance Decay introduced compaction-induced governance decay --- the exact phenomenon studied here --- along with the compaction-eviction attack this paper's threat model adopts, the constraint-pinning defense (a compaction-exempt pinned buffer with verbatim re-injection and per-step integrity checking) that three of this paper's five variants decompose, its defender-side cost ($\approx$47 tokens, $<$0.5\% overhead), and the finding that an optimized injection defeats pinning on every tested model, with operator impersonation in recent context --- the ancestor of this paper's template-role-forgery family --- as the residual failure mode. What it does not report is any of this paper's estimands: its result is binary defended/undefended at fixed attack configurations, with no continuous strength axis, no break-boundary location, no variant-resistance ordering, and no attacker-to-defender cost ratio. Two further neighbors bound specific instruments. IH-Benchmark (arXiv 2607.25987) supplies the nearest oracle design --- a predicate DSL plus category-scoped LLM judges over instruction-hierarchy conflicts --- which this paper's two-tier oracle follows structurally; the delta is confined to the before/after-compaction persistence gate. Security--Fidelity Tradeoffs (arXiv 2606.30783) prices the hidden fidelity cost of prompt-injection defenses --- the defender half of this paper's cost account --- without an attacker-side denominator or a boundary to evaluate it at.

Among earlier evaluation designs and mechanisms, TensorTrust builds an interpretable, human-authored attack/defense benchmark and explicitly measures a robustness-vs-usability trade-off --- the closest existing evaluation design to what this paper needs --- but its defenses are static and single-turn, with no compaction dynamic and no continuous attack-strength axis to sweep. LightMem is the closest existing mechanism: its ``sleep-time'' consolidation step is structurally similar to the protected-memory slot this paper's structural-slot variant is modeled on, but LightMem is evaluated only on benign long-horizon dialogue consistency, never under adversarial stress --- precisely the gap this paper's adversarial sweep is built to close.

\textbf{Prompt-injection attacks and defenses for LLM agents.} ChatInject wraps malicious tool-output content in forged chat-role delimiters to make an agent treat attacker text as higher-privilege structure, evaluated for attack success at one fixed configuration against prompt- and detector-based defenses. ThreatSnapshots isolates a single backbone-LLM call with its full execution context to score and compare backbone security across models without simulating whole agents. Both establish that LLM-agent security can be measured and compared, but neither sweeps a continuous attack-strength parameter against structured memory-pinning defenses, and ThreatSnapshots's single-call abstraction cannot represent a pinned constraint's survival, which is a property of a repeated, multi-event lifecycle rather than any one call.

\textbf{Optimization-based jailbreaks.} GCG established gradient-guided greedy coordinate-descent search over adversarial suffix tokens to elicit an affirmative response from an aligned model; SlotGCG extends this by making insertion position itself an optimized variable, and TAO-Attack replaces GCG's single prefix-target loss with a progressive two-stage objective. All three optimize against single-turn safety-refusal at a fixed attack budget, not a swept continuous strength variable against a structured pinning defense, and none measures the defender-side cost the attack must overcome --- leaving the attack/defense cost trade-off under compaction unmeasured, which this paper's cost-accounting method is built to close.

\textbf{Injection-success evaluation and robustness benchmarking.} GuidedBench sharpens jailbreak-success evaluation into an entity/action checklist scored by a guideline-based judge, reducing evaluator ambiguity --- but its object is single-turn harmful-content elicitation, providing no oracle for whether a pinned constraint was violated after a compaction event, which requires checking persistence across time rather than harmfulness of one response. IH-Benchmark, discussed above, comes closest on the instrument side --- its predicate-DSL-plus-scoped-judge design is the structural template this paper's oracle follows --- but scores single-decision conflict outcomes with no persistence gate. TensorTrust, discussed above as the nearest earlier evaluation design, measures robustness and task validity jointly in a static, single-turn game.

\textbf{Context compression and memory pipelines.} COMI frames long-context compression as constrained information selection under a fixed token budget, allocating that budget by query relevance and redundancy; a companion information-theoretic view models a compressor as a noisy channel and measures compression quality by mutual information between source and compressed representation. Both model what compression cooperatively retains under a benign objective; neither models an adversary actively trying to evict a protected constraint, so their retention guarantees say nothing about adversarial durability --- the gap this paper's RQ0 attribution check directly verifies. LightMem, discussed above as the nearest mechanism neighbor, restructures the memory lifecycle into pre-compression, topic segmentation, summarization, and offline consolidation, giving a concrete, realistic pipeline for preserving persistent content across sessions, but stress-tests none of it adversarially.

\section{Conclusion}

This work reframes a governance assumption behind memory-pinning --- that a constraint fixed once survives every later compaction pass --- as a precisely measurable boundary and price, rather than a binary robust/not-robust property. The central methodological wager is that pinning's failure is only informative once it is located along a comparable strength axis and priced in a common attacker/defender currency; a bare ``it breaks'' finding, already established by Governance Decay --- whose pinning defense the variants studied here decompose --- tells a deployment decision nothing about which variant to choose or when to worry. Whichever ordering and cost differentiation the sweep reveals, the finding is decision-relevant: a cheap, low-strength boundary for the cheapest variant would argue against relying on pinning alone at any meaningful adversary budget, while a well-separated, expensive-to-reach boundary for the structural-slot variant would argue that mechanism choice, not just ``pinning versus no pinning,'' is where the real safety margin lives. Three limitations bound this claim from the outset and are natural next steps rather than afterthoughts: the cost ratio's currency-dependence is treated as a first-class, reportable result rather than resolved by picking one currency, and a genuine reversal across currencies would itself be a finding, not a design flaw to hide; the warm-started, transferable gradient attacks used to keep the white-box arm computationally tractable may understate true attack strength relative to a fully re-optimized attack, a calibration gap this design measures on a subset rather than assumes away; and the gradient family's need for open weights confines that arm to open-weight backbones, so cross-model conclusions for closed-weight API models rest on the black-box template family alone.

\vspace{1em}
\hrule
\vspace{1em}

\section*{Pre-handoff checklist}
\begin{itemize}[leftmargin=1.4em]
\item \textbf{Organization/navigation} --- narrative flows problem $\to$ method $\to$ experiments $\to$ related work without forward references to undefined terms.
\item \textbf{Terminology discipline} --- ``pinned constraint,'' ``compaction event,'' ``break-through,'' ``boundary,'' ``cost ratio'' each defined at first use in Problem Definition, used consistently thereafter; no internal stage IDs (\texttt{sp\_*}, \texttt{m\_*}, \texttt{rq\_*}) leaked into prose --- checked directly against the discussion pipeline log.
\item \textbf{Notation with worked example} --- the formal setup section defines $(T, \kappa, c, C, d, D, a, s, V, p_{d,a,g})$ and the Problem Definition/Methodology sections consistently reuse this notation; no separate worked numeric example is included since no numbers exist yet (would require fabricating results).
\item \textbf{RQ framing matches outcomes} --- all 6 RQs are graded/quantitative (rates, boundary values, rank-correlations, agreement coefficients), each answered in its own subsection with an explicit falsifiable expected outcome and a stated falsification condition, not a binary yes/no.
\item \textbf{No typesetting artifacts} --- only explicit placeholder markers in Results, per the drafting convention this is a pre-registered spec, not a report.
\item \textbf{Artifact plan} --- the defense-variant middleware, the strength-normalization scheme, the two-tier oracle, and the constructed pinned-constraint suite are the natural release artifacts; not yet built, so not claimed as available.
\end{itemize}

\section*{Hand-off}
This is a draft, not a finished paper --- experiments have not run. Natural next steps: prose polish and venue-specific LaTeX formatting once a target venue is chosen, or revising an earlier stage if piloting surfaces a problem (e.g., if the oracle's judge-human $\kappa$ in RQ3 is much lower than expected).

\end{document}
